% Source: https://artofproblemsolving.com/community/c3h3180052_2022_lmt_spring_speed_round__lexington_mathematical_tournament
% Stable identity: aops-topic-3180052-q19
If a parent side is $s$, a child side has perpendicular components $\dfrac{3s}{7},\dfrac{4s}{7}$, so its squared length is $\dfrac{25}{49}s^2$. Thus areas form a geometric series with ratio $\dfrac{25}{49}$. Their sum is $\dfrac{1}{1-\dfrac{25}{49}}=\dfrac{49}{24}$, giving $\boxed{a+b=73}$.
\begin{asy}
usepackage("amsmath");
size(7cm);
pair[] p={(0,0),(1,0),(1,1),(0,1)};for(int k=0;k<5;++k){path q;for(int i=0;i<4;++i)q=q--p[i];draw(q--cycle);pair[] v;for(int i=0;i<4;++i)v.push((4*p[i]+3*p[(i+1)%4])/7);p=v;}label("successive ratio 3:4 divisions",(.5,-.2));
\end{asy}
